\documentclass[11pt,a4paper]{article}
\usepackage[T1]{fontenc}\usepackage{lmodern,amsmath,amssymb,amsthm,mathtools,booktabs,microtype}
\usepackage[margin=27mm]{geometry}\usepackage[colorlinks=true,linkcolor=blue,urlcolor=blue,citecolor=blue]{hyperref}
\newtheorem{theorem}{Theorem}[section]\newtheorem{lemma}[theorem]{Lemma}\newtheorem{proposition}[theorem]{Proposition}
\newcommand{\Q}{\mathbb Q}\newcommand{\Z}{\mathbb Z}\DeclareMathOperator{\Tr}{Tr}
\setlength{\emergencystretch}{2em}
\title{A quartic in 204 variables from recovered elliptic differences}
\author{}\date{14 November 2026}
\begin{document}\maketitle
\begin{abstract}
Under the stated inherited arithmetic premises, we construct a family
of integer quartics in 204 total variables, eight fewer than the
212-variable predecessor. The family is a sum of 186 primitive
quadratic squares and has exactly 3046 ordinary monomials. Two elliptic
differences are recovered from reciprocal point pairs instead of being
stored separately. A positive weighted trace controls the enlarged
point domain, including half-integral indices and torsion differences.
The old Gaussian congruences also permit the required half-generator
in the completeness argument. The coefficient-height improvement is
preserved. The sufficient defining parameter remains ineffective.
\end{abstract}

\section{Statement and precise scope}
All variable counts include the input. Set
\[
 K=\Q(\rho),\quad \rho^2=2,\quad d=3-\rho,\quad D=d^2,
 \quad E:y^2=x^3-Dx,
\]
\[
 R_0=(1+\rho,2\rho),\qquad P=2R_0.
\]
We retain $2E(K)=\Z P$ with $P$ non-torsion, and the precise
scalar, pole, contact-height, and fixed-data prime-selection inputs
of the 212-variable source~\cite{Parent}. These are implications with
explicit arithmetic premises; the finite coefficient checks do not
prove them. The full external arithmetic audit is not declared complete.
No additional $abc$ or Szpiro hypothesis is introduced.

\begin{theorem}\label{thm:main}
Under these inputs, there is a computable family
\[
 P_N\in\Z[X,Y_1,\ldots,Y_{203}],\qquad N\ge1,
\]
with degree exactly four, integer content one, and exactly 3046 nonzero
coordinate monomials. There is one fixed integer $N_0$ such that every
$N\ge N_0$ has rational input projection exactly $\Z$.
No numerical sufficient $N_0$ is determined.

For the balanced presentation put $C=149^N$, $K_0=31581660^{28}$,
$J_0=5898050446009651968$, and $H_0=(J_0K_0)^2$. Its exact coefficient
height, for every $N\ge1$, is
\[
                    H(P_N)=\max(H_0,K_0C^2).
\]
It equals $31581660^{28}149^{2N}$ for $N\ge57$. This is a coefficient
threshold, not a sufficient defining threshold.
\end{theorem}

The complete quadratic and expanded quartic tables, their producer,
and the independent reconstruction checker are in
\begin{center}\small
\nolinkurl{computation/quartic/reciprocal204_v27}.
\end{center}
They specify every coefficient, not just a variable count or an
unexpanded expression. The arithmetic proof is separate below.

\section{Retained equations and arithmetic conclusions}
Start with the unbalanced 212-coordinate table, SHA-256
\begin{center}\small
\texttt{b2f32d64ef1fb2591d37cfd41aadf8869ce691c39ed92e0e597ec10aaa53fc75}.
\end{center}
All indices in the transformation below refer to that table before
compaction. Write $R=x_{118}$ and $R_2=x_{117}$. The retained scalar
rows, reference-point rows and second root block imply
\begin{equation}\label{eq:scalar-identities}
 \begin{gathered}
 \beta=51X^2+2504,\quad
 r=\beta(\beta-1)(\beta-2401)(32\beta+343)>0,\\
 (1+149^2r)R=149^{2N+1}r,\qquad R_2=R^2.
 \end{gathered}
\end{equation}
In particular $R>149^3/(1+149^2)>100$. For the strict lower bound,
$r\ge1$ on the rational input domain, since $\beta\ge2504$.
The retained scalar point is $Q_2=[a]P$ with $a\ne0$. Its proof uses
the reference point, this scalar block and the scalar pole, not the
first or third root block or the aggregate equation being replaced.
The same proof gives
\begin{equation}\label{eq:scalar-bound}
 c=|a|\ge149^{L_N},\quad
 h(\beta)\le C_3\,149^{-L_N}c(1+\log c/3),\quad
 L_N=\left\lfloor\frac{9(2N-1)}4\right\rfloor+1,
\end{equation}
where $C_3=2^{330770}$. Its uniform completeness statement has one
fixed scalar threshold, before the later input-dependent prime choices.
These are the exact source scalar assertions retained here.

The fourth root block is also retained: $Q_4=[b]P$, $b\in\Z\setminus\{0\}$.
Nonvanishing of $b$ is essential in the Lucas argument; this block is
not removed. The reference point $H$, the point $H-P$, and their
parameters
\[
 S=y(H)/x(H),\qquad T_H=y(H-P)/x(H-P)
\]
are recovered by the unchanged reference equations and scalar argument.
They are non-torsion and their indices avoid 0 and 1, as in the source.
Neither reference ordinate parameter is zero.

The contact set and exceptional primes remain
\[
 \Theta=\Theta_0\cup(2-\Theta_0),\qquad
 \Theta_0=\{0,-1/2,-1,1/2399,-32/407\},
\]
\[
 \Sigma=\{p:p\le76176\}\cup\{2029147\}.
\]
For $t\notin\Theta$, let $M(t)$ be the squarefree product of primes
outside $\Sigma$ for which one $v_p(t-\theta)$ is positive and odd.
The same full-family modular-height input gives
\begin{equation}\label{eq:heightcontact}
                 h(t)\le H_2 M(t)^{7/6}
\end{equation}
with a fixed constant, including $t=1$, for which $M(t)=1$.
The proof applies the inherited five-contact bound to both reflected
blocks; no new contact family is used here.

\section{Reciprocal point pairs, including every exceptional case}
For $Q\in E\setminus\{(0,0)\}$ write
\[
 f(Q)=1/x(Q),\qquad h(Q)=-y(Q)/x(Q)^2,
\]
extended by $f(O)=h(O)=0$. The reciprocal equation is
$h^2=f-Df^3$. Its coordinates are ordinary affine coordinates after
translation by the two-torsion point $(0,0)$, via $(f,h)\mapsto(-Df,-Dh)$.

For each $i=1,3$, retain five field coordinates $f_i,h_i,g_i,k_i,s_i$
and use the existing reciprocal-square outputs $F_i,G_i$. Impose
\begin{equation}\label{eq:pair}
 \begin{split}
 F_i=f_i^2,\quad G_i=g_i^2,\qquad h_i^2=f_i-Df_iF_i,\\
 s_i(f_i-g_i)=h_i+k_i,\qquad
 s_i(h_i-k_i)=1-D(F_i+f_ig_i+G_i).
 \end{split}
\end{equation}
These are five field quadrics. Put, without adjoining outputs,
\begin{equation}\label{eq:diff}
 x_i=s_i^2+D(f_i+g_i),\qquad
 y_i=Dh_i-s_i(Df_i+x_i).
\end{equation}

\begin{lemma}\label{lem:pairs}
Every solution of \eqref{eq:pair} describes an ordered pair
$U_i,V_i\in E\setminus\{(0,0)\}$, $U_i\ne V_i$, and their finite
difference $Q_i=U_i-V_i=(x_i,y_i)$. Conversely every such ordered
pair supplies a unique $s_i$. Identities, tangent pairs and finite
torsion differences are included.
\end{lemma}
\begin{proof}
Suppress the index and put
$r_1=s(f-g)-h-k$, $r_2=s(h-k)-1+D(f^2+fg+g^2)$,
$r_3=h^2-f+Df^3$. Direct polynomial identities are
\[
 k^2-g+Dg^3=r_3-(f-g)r_2+(h-k)r_1,
\]
\[
 y^2-x^3+Dx=D^2r_3-D(Df+x)(r_2-sr_1).
\]
They put both reciprocal points and the asserted difference on the
curve. If $f\ne g$, the secant through the first translated point
and the negative of the second has slope $s$, and the chord law gives
\eqref{eq:diff}. If $f=g$, then $k=-h$. The case $h=0$ would imply
$f=0$ or $Df^2=1$, whereas the last equation demands $1-3Df^2=0$;
this is impossible. Thus $h\ne0$ and $2sh=1-3Df^2$, exactly the
tangent slope, giving the same difference formula. Conversely the
chord or tangent law gives the displayed equations. The reciprocal
origin $(0,0)$ represents $O$ on the original curve, so no identity
case has been divided out.
\end{proof}

The two ratios for each pair are specified by four additional field
quadrics, with $\kappa_1=1$, $\kappa_3=2$:
\begin{equation}\label{eq:ratios}
 \begin{split}
 h_iL_i+\kappa_i f_iS=0,&\quad
 (1-DF_i)L_i+\kappa_i h_iS=0,\\
 k_iL_i^0+\kappa_i g_iT_H=0,&\quad
 (1-DG_i)L_i^0+\kappa_i k_iT_H=0.
 \end{split}
\end{equation}
On the nonidentity domain these give
$L_i=\kappa_i z(U_i)/z(H)$ and
$L_i'=\sigma L_i^0=\kappa_i\sigma(z(V_i)/z(H-P))$, with $z=-x/y$.
At $U_i=O$ the second equation forces $L_i=0$, and similarly for $V_i$.
A nonidentity two-torsion numerator has $h_i=0,f_i\ne0$ and would
force $S=0$, contrary to the retained reference domain. Thus no such
numerator supplies a spurious ratio. The difference itself may be
finite two-torsion; it is not excluded.

Write
\[
 \nu(Q_1)=n\in\tfrac12\Z,\quad Q_2=[a]P,\quad
 2\nu(Q_3)=t\in\Z,\quad Q_4=[b]P.
\]
The factor two in \eqref{eq:ratios} makes the third recovered label
$t$, not $t/2$. This preserves the integral Lucas equation.

\section{Exact transformation of the stored rows}
Delete the two seven-field root blocks, old coordinates 155--168 and
183--196, and their respective 18 rational rows 104--121 and 140--157.
Append the ten new field coordinates in \eqref{eq:pair}, with the
nine field quadrics \eqref{eq:pair}--\eqref{eq:ratios} for each pair.
The count is $212-28+20=204$ rational coordinates and
$186-36+36=186$ rational rows.

The reused reciprocal-square slots are $(54,55),(23,24)$ for pair 1,
and $(58,59),(27,28)$ for pair 3. Their second coordinate is one half
of the $\rho$ coefficient. Ratio slots are $(31,32),(33,34)$ and
$(39,40),(41,42)$, respectively. The reference slots are
$(19,20),(211,212)$ with ordinary component conventions.

Before compaction make the invertible changes
$x_{i,\mathrm{old}}=R^{e_i}x_{i,\mathrm{new}}$, where
\[
 e_i=1\ (39\le i\le46\text{ or }48\le i\le51),\qquad
 e_i=-1\ (205\le i\le208),
\]
and all other $e_i$ vanish. In each affected row collect the total
power of $R$, replace $R^2$ by its existing slot $R_2$, and divide
by its least occurring power of $R$. Since $R>0$, this preserves its
zero set. Every remaining power is 0, 1 or 2 and the resulting row
is quadratic. This is checked literally for all 22 affected rows;
negative least powers are allowed. It unscales the third and fourth
label ratios and the two Lucas products, while adjusting the fourth
root ordinates inversely. It neither changes the scalar block nor
assumes a numerator nonzero.

Replace the old aggregate row 180 by the single rational quadratic
\begin{equation}\label{eq:aggregate}
 8(\lambda-t_1)+4\operatorname{Re}(x_1+x_3)
 +\operatorname{Re}\bigl((C_2-E_2)^2+(C_4-E_4)^2\bigr)=0.
\end{equation}
Here $C_2,E_2$ are old field slots $(173,174),(175,176)$ and
$C_4,E_4$ are $(201,202),(203,204)$. The retained root identities are
$x(Q_j)=(C_j-E_j)^2/(4R^2)$ for $j=2,4$. Consequently
\[
 t_1=\lambda+T,\qquad
 T=\tfrac14\Tr\bigl(x(Q_1)+x(Q_3)+R^2(x(Q_2)+x(Q_4))\bigr).
\]
By \cite[Corollary~4.1]{Weighted}, $T>0$, and, for
$B=\max(|n|,|a|,|t|,|b|)$, there are fixed constants
\begin{equation}\label{eq:weighted}
                         B^2\le A(h(R)+h(T))+J.
\end{equation}
The weighted height theorem is proved for precisely these enlarged
point domains. Reusing the old unweighted theorem here would not be
justified.

Keep all contact, local-factor, trace, and pole equations, then compact
coordinates in their old increasing order and append the new pairs.
Clear every rational row to a primitive integral row with the specified
sign convention. Their sum of squares is the unbalanced quartic.
The old storage of $(L_1-\beta+3)^2$, rather than $(L_1-\beta)^2$,
is retained: its later fourth-power equation performs the correct
shift. The constant 2501 in that intermediate row is intentional.

\section{Local extraction with half-integral indices}
The unchanged local factors, writing $q$ for the contact weight, are
\begin{align*}
1+q^5C_2^2+q^2D_H^2+q^4B_1^4,&\qquad
1+q^5G^2+q^2A_3^2+q^4B_0^4,\\
1+q^5I^4+q^3A_2^2+q^2A_4^2+q^4B_3^4,&\\
1+q^5C_3^2+q^2C_1^4+q^4B_4^4,&\qquad
1+q^5C_4^2+q^2A_1^2+q^4B_2^4.
\end{align*}
Here $D_H=f(H)$, $A_i=f(U_i)$, $B_i=\sigma f(V_i)$,
$B_0=\sigma f(H-P)$, $C_i=L_i-L_i'$, $I=L_1-\beta$, and
\[
 G=L_3^2-L_2(L_2-2)L_3L_4+(L_2-2)^2L_4^2-1.
\]
The factor conventions, including stored affine shifts, are those
verified by the exact row transformation.

At a consumed prime, $v_p(q)=-\mu$ with $\mu>0$ odd. Every finite
valuation of a reciprocal abscissa is even: this follows directly
from $h^2=f-Df^3$ when $v(f)\ne0$, and valuation zero is even too.
Zeros may be omitted. The distinct valuation classes in the five
factors therefore remain the same as in the source, even for arbitrary
points, half-integral indices and zero numerator points. The pole
obstruction and reciprocal trace give, at one place above each prime,
\[
 v(C_2,C_3,C_4,G)\ge3,\quad v(C_1,I)\ge1,
 \quad v(D_H,A_i,B_i)\ge2.
\]
This uses arbitrary finite test valuations, not only positive ones.

For clarity, the formal-group step is also checked on the enlarged
domain. Let $\nu(H)=h=m/2$ and $\nu(U)=u\in\tfrac12\Z$.
The identities
\[
 [2m]U=[4u]H,\qquad [2(m-2)]V=[4(u-\nu(Q))](H-P)
\]
kill all two-torsion and give exact logarithmic ratios
$u/h$ and $(u-\nu(Q))/(h-1)$. At least one of $h,h-1$ is a
$p$-adic unit since $p$ is odd. One ratio is therefore integral.
The positive comparison test transfers integrality to the other
approximate ratio and then to its exact ratio. Only after this step
may the relative formal-logarithm error of order four be used as an
absolute congruence. For the third pair multiply both ratios by 2,
a unit at these primes. Subtracting the two congruences gives
$L_1\equiv n\pmod p$ and $L_3\equiv t\pmod {p^3}$.
The unchanged second and fourth blocks give $a,b$ modulo $p^3$.
Thus for both arguments $t_0=\lambda,t_1=\lambda+T$,
\begin{equation}\label{eq:divisibilities}
 M(t_j)^3\mid\Delta,\qquad
 M(t_j)\mid J_X\quad\text{if }J_X\ne0,
\end{equation}
where
\[
 \Delta=t^2-a(a-2)tb+(a-2)^2b^2-1\in\Z
\]
and $J_X$ is the positive absolute reduced numerator of $\beta-n$.
Identity numerator points are covered by $z(O)=0$ and the second
ratio equation. The same argument includes inert places; incompatible
formal requirements simply exclude a solution there.

\section{One uniform soundness threshold}
Elementary height estimates applied to \eqref{eq:scalar-identities}
give fixed constants
\begin{equation}\label{eq:Rheight}
                     h(R)\le4h(\beta)+C_RN+C_0.
\end{equation}
Indeed the scalar contact polynomial has degree four and its
fractional-linear normalization adds at most the height of
$149^{2N+1}$ and a fixed constant. From
$c\ge149^{L_N}$, $c\le B$, it follows that $N\le C\log B+C'$.
Also \eqref{eq:scalar-bound} implies the uniform coarse estimate
$h(\beta)=O(B(1+\log B))$, independently of $N$ and the solution.

If $M(t_j)^k\mid J_0\ne0$ separately for the two arguments, then
\eqref{eq:heightcontact}, $h(T)\le h(t_0)+h(t_1)+\log2$, and
\eqref{eq:weighted}--\eqref{eq:Rheight} imply
\begin{equation}\label{eq:transfer}
 B^2\le C|J_0|^{7/(6k)}+C'h(\beta)+C''N+C'''.
\end{equation}
All constants are fixed before $N,X$ and the witnesses.

If $\Delta\ne0$, then $|\Delta|\le14B^4$ for $B\ge1$.
Take $k=3$ in \eqref{eq:transfer} to obtain
\[
 B^2\le C_1B^{14/9}+C_2B(1+\log B)+C_3'\log B+C_4.
\]
Every term on the right is $o(B^2)$, so $B$ is bounded by a fixed
constant. Increasing a single $N$ threshold makes $c\ge149^{L_N}$
exceed that bound. Hence $\Delta=0$ for all solutions at all larger $N$.

The retained $b$ is nonzero. If $t=0$, the equation $\Delta=0$
forces $|(a-2)b|=1$, impossible for $c=|a|$ sufficiently large.
Thus $a,t,b$ are all nonzero. The source elementary Lucas argument
now applies with integer labels: writing $w=\operatorname{sgn}(a)(a-2)b$
gives $t^2-ctw+w^2=1$. Its positive norm-one solutions are consecutive
terms of $U_0=0,U_1=1,U_{j+1}=cU_j-U_{j-1}$ up to signs and
interchange. The divisibility by $a-2$, reduced modulo $c-2$ or $c+2$,
forces the relevant index at least $c-2$. The same source bound is
\begin{equation}\label{eq:Lucas}
                   \log B\ge(c-5)\log(c-2)\qquad(c\ge8).
\end{equation}
This is why the fourth block, including $b\ne0$, was retained.

Suppose now $J_X\ne0$. Writing $\beta=u/v$ reduced and $n=l/2$
gives
$|J_X|\le|2u-lv|\le4B\exp h(\beta)$, since $|n|\le B$ and
$B\ge1$. The extra height terms in \eqref{eq:transfer} are at most
$B^2/2$ once $B$ exceeds a fixed constant, by the coarse uniform
bounds above. The forced lower bound on $c$ ensures this at all
sufficiently large $N$. With $k=1$ we obtain
\[
 B^{5/6}\le C_5\exp(7h(\beta)/6),\qquad
 \log B\le\tfrac75h(\beta)+C_6.
\]
Choose $N$ large enough that $(7/5)C_3 149^{-L_N}\le1$.
This upper bound is at most
$c+(c/3)\log c+C_6$, contradicting \eqref{eq:Lucas} for all
sufficiently large $c$. Hence $\beta=n\in\tfrac12\Z$.
For $X=u/v$ reduced, $v>0$, this forces $v^2\mid102$.
As 102 is squarefree, $v=1$. Soundness holds with one fixed
ineffective threshold, not an input-dependent choice of $N$.

\section{Completeness and the half-generator}
Fix an integer input and any $N$ beyond the source's fixed scalar
completeness threshold. The unchanged scalar preparation gives a
positive integer $a\ge3$, its point $Q_2=[a]P$, the normalizer $R$,
and a reference progression $h\equiv h_*\pmod L$ with
$64\mid h_*,L$, preserving all scalar and finite-place conditions.
Choose positive Lucas labels $t,b$ with $\Delta=0$ as in the source:
for $j>0$ divisible by $a-2$, take $b=U_j/(a-2)$ and $t=U_{j+1}$.
Set
\[
 Q_1=[\beta]P,\quad Q_3=[t]R_0,\quad Q_4=[b]P.
\]
These points are nonzero. The weight $R$ and the new $T>0$ are now
fixed, before any comparison prime is selected.

For an integer reference $h$, the required points are
\[
 H=[h]P,\quad U_1=[h\beta]P,\quad V_1=[(h-1)\beta]P,
\]
\[
 U_3=[ht]R_0,\quad V_3=[(h-1)t]R_0,
\]
and the source points for blocks 2 and 4. Their differences are the
chosen $Q_i$. Compact real neighborhoods can be chosen avoiding the
finite chart exceptions. There is no rational multiplication map
by $1/2$ being used: with $h=64k$, all displayed points are fixed
integer multiples and translates of $[k]R_0$. The source compact-group
recurrence, after enlarging the component modulus, therefore supplies
joint neighborhoods in this fixed product group. These neighborhoods
may depend on the fixed labels, but impose no new lower bound on $N$.

Use the additive contact coefficient list
$\Theta\cup(\Theta-T)$, merge repetitions while retaining their
incidences, and clear its finitely many rational denominators by the
source square dilation. All coefficients, forbidden primes, congruence
moduli and real neighborhoods are fixed before applying the inherited
fixed-data simultaneous-prime theorem. It applies to this arbitrary
finite rational list, exactly as in the 212 proof.

The selected Gaussian generators already have real part $9\pmod{16}$
and imaginary part $0\pmod{16}$. By \cite[Lemma~5.1]{Weighted}, their
old annihilating moduli $\#E(\mathbf F_p)/4$ kill $R_0$ itself in
both conjugate groups. Therefore the same relatively prime moduli
$U,V$ permit $h\equiv0\pmod U$, $h\equiv1\pmod V$ and
$h\equiv h_*\pmod L$. These congruences make the new third pair
formal at its designated places, just as for the doubled generator.
This is the additional completeness check required by replacing a
$tP$ root by a $tR_0$ root.

At the selected places the full ratios, with multipliers $1,2$ on
the new pairs, are the exact labels $\beta,t$. All comparison errors
have valuation at least four, and reciprocal abscissas have valuation
at least two. The unchanged factor exponents consequently give unit
factors at the chosen place, and reciprocal-trace residue one or two
at the other. At other primes all nonselected contact factors
contribute square classes as before. At the prepared finite primes
the weight tends to zero; at the two real embeddings it is positive
and tends to zero, uniformly on the fixed compact neighborhoods.
Thus the actual pole inputs satisfy the source square-denominator,
local squareclass, depth and real interval requirements. The retained
pole-realization theorem supplies their independent witnesses.

Finally use the reciprocal secant/tangent formulas, the ratio
identities, the inverse rescalings by $R\ne0$, and the retained
contact circuits to fill every new coordinate. They satisfy all 186
quadrics. This is a construction of fresh witnesses, not an asserted
bijection with every old witness. All input-dependent choices occur
after a fixed scalar threshold. Combining this with the fixed
soundness threshold proves the quantifier order in Theorem~\ref{thm:main}.

\section{Coefficient balancing and complete checks}
In the preceding unbalanced rows $A=149^{2N+1}$. Put $C=149^N$ and
substitute $A=149C^2$. Rescale the two retained slots $R,R_2$ by $C$;
in each row divide by the least resulting power of $C$, then make the
integer row primitive. Since $C\ne0$, this is an invertible rational
change at every fixed $N$. All parameter degrees after division are
at most one and all coordinate degrees at most two.

The checker independently reconstructs every deleted, changed and
retained row from the pinned 212 source, including coordinate
conventions and the two ratio multipliers. It expands each row square
by an unordered term-pair sum, rather than importing the producer's
expansion. Both presentations have 3046 coordinate monomials; the
unbalanced and balanced tables have 3047 and 3048 terms when fixed
parameter powers are recorded separately. Every coordinate occurs.
An actual fourth-power coefficient of one proves degree four and
integer content one.

For the balanced quartic, all coefficient profiles are either linear
$c_1C$ or of the form $c_0+c_2C^2$. Exact integers in the checker prove
\[
 c_1^2\le H_0K_0,
 \qquad |c_0|K_0+|c_2|H_0\le H_0K_0,
\]
respectively. Profiles attaining $H_0$ and $K_0C^2$ occur. These
inequalities imply the exact maximum in Theorem~\ref{thm:main} for
every positive $C$. The exact comparison
$K_0 149^{112}<H_0\le K_0 149^{114}$ proves the threshold 57.
Four deliberately corrupted row, ratio, aggregate and quartic data
sets are rejected. These are finite algebraic certifications, not a
formal proof of the inherited arithmetic statements or an independent
non-author review.

\begin{thebibliography}{9}\small
\bibitem{Parent} \emph{A 212-variable quartic with additive height comparisons},
project companion, 13 November 2026,
\nolinkurl{quartic212_20261113_v25.tex}; exact source tables in
\nolinkurl{computation/quartic/additive212_v25}.
\bibitem{Weighted} \emph{Heights of positive weighted elliptic traces},
project companion, 14 November 2026,
\nolinkurl{weighted_trace_20261114_v27.tex}.
\bibitem{Reciprocal} \emph{Trace heights on quadratic twists with $j=1728$},
project companion, 13 November 2026,
\nolinkurl{trace_on_all_points_20261113_v25.tex}, Section~2.
\bibitem{Source} Pinned 214-coordinate source and arithmetic applications,
\texttt{elliotglazer/diophantine-z-q}, commit
\texttt{1f7e1f1fb0438feda3c9f3da9cc16982126aeeed},
\nolinkurl{research/quartic214/proofs/uniform-existence.md}, P03--P20.
Numeric source and provenance retained in
\nolinkurl{computation/quartic/balanced214_v24/source}.
\end{thebibliography}
\end{document}
